\documentclass[11pt]{article}
\usepackage[a4paper,margin=2.5cm]{geometry}
\usepackage{graphicx}
\usepackage{float}
\usepackage{booktabs}
\usepackage{amsmath}
\usepackage[hidelinks]{hyperref}
\setlength{\parindent}{0pt}
\setlength{\parskip}{6pt}

\newenvironment{reviewquote}%
  {\begin{quote}\itshape}%
  {\end{quote}}

\begin{document}

\begin{center}
{\Large\bfseries Response to Reviewers}\\[4pt]
{\large Comparing One-Layer and Two-Layer Economic MPC\\
for a SAG Milling Circuit -- AdCONIP 2026}
\end{center}

We thank both reviewers for their careful reading and constructive
comments. Each comment is quoted in italics, followed by our response.
All changes described below are included in the revised manuscript.

\section*{Reviewer 1}

\subsection*{Comment 1 -- Typo in Eq.~(2)}

\begin{reviewquote}
There is a typo in Eq.~(2).
\end{reviewquote}

\textbf{Reply.} Thank you for catching this. It is a font-substitution
artifact in the Word version of the manuscript. The abbreviation
``s.t.'' (subject to) in Eq.~(2) was stored as a text run inside the
equation. The latex and the PDF files are correct. We have rewritten the
abbreviation as roman math letters. Reviewer 2 reported the same issue
(Minor comment 2), so both reports are resolved by this fix.

\section*{Reviewer 2}

\subsection*{Comment 1 -- Symbol inconsistency for mill speed}

\begin{reviewquote}
The mill speed is labelled $\omega_c$ in Fig.~1 but $\varphi_c$ in
Section VII-A. Since $\varphi_f$ and $\varphi_c$ would then coexist
with different meanings, this risks confusion. The variable should be
defined once in Section II and used consistently throughout, including
figures.
\end{reviewquote}

\textbf{Reply.} We checked Fig.~1 at full resolution: the mill speed is
labelled $\phi_c$ there, and $\omega_c$ does not appear anywhere in the
manuscript or its figures. The reviewer likely read a low-resolution
rendering, where the small italic $\phi$ can be mistaken for an
$\omega$. The underlying request is implemented together with
Comment~2: Section~II now presents the main model equations and defines
the raised variables at first use, $\phi_c$ (mill speed) and $\phi_f$
(ore hardness). The reviewer's concern is on point, the two symbols are
easy to confuse. However, we keep the $\phi$ notation of the underlying
model literature, and all variable descriptions are included in the
manuscript.

\subsection*{Comment 2 -- Section II lacks model equations and parameter definitions}

\begin{reviewquote}
The key state equations, power-draw expression, and fines-production
relation should be included. The parameters $\alpha_r$ and $\varphi_f$
should be introduced here as part of the model description.
\end{reviewquote}

\textbf{Reply.} Agreed, and implemented. Section~II now presents the
main governing equations of the model, the state balances, the breakage
relations and the power draw, and defines the main variables at first
use, including $\alpha_r$, $\phi_f$ and $\phi_c$. The parameters
$\alpha_r$ and $\phi_f$ are thereby introduced before their use in
Section~V.

\subsection*{Comment 3 -- The assumption $A=\mathrm{const}$ lacks quantitative justification}

\begin{reviewquote}
The reader needs at least an order-of-magnitude estimate of how much
$A$ varies over the PSE range used (0.65 to 0.715).
\end{reviewquote}

\textbf{Reply.} This comment is right on the spot. Historically, we
used a step function for the PSE floor, with two constant values (0.65
and 0.715), to test how each control system behaves when the set-point
changes. We also did not simulate the downstream circuits, flotation
and leaching, which are the parts that actually compute and use
$A(\mathrm{PSE})$.

In our case, however, the choice of one or two set-points does not
change the result. The study is an apples-to-apples comparison of the
control designs under one common protocol. As long as the closed-loop
PSE hovers close to the active set-point, $A(\mathrm{PSE})$ can be
treated as a constant to simplify the optimisation. Figure~\ref{fig:pse}
shows the realised PSE of all six benchmark runs against the stepped
floor: within each phase of the protocol the closed-loop PSE stays
within a few hundredths of the active floor.

\begin{figure}[H]
\centering
\includegraphics[width=\textwidth]{figs/pse_vs_floor_both_objectives.png}
\caption{The v5 ore-hardness realisation (top) and the raw closed-loop
PSE of the three architectures under the energy and the profit
objective against the stepped product-size floor (0.65 to 0.715 at
75~h). Within each phase the PSE hovers at the active floor, which
justifies treating $A(\mathrm{PSE})$ as constant in the benchmark.}
\label{fig:pse}
\end{figure}

We also recomputed SEC using only the region after the 75~h step, where
a single floor value is active (Table~\ref{tab:sec75}). The result
becomes stronger: the SEC values cluster even more tightly than over
the full protocol, so the single-set-point view reinforces the paper's
conclusion that specific energy is comparable across the
architectures. The common upward shift of about 0.3~kWh/t reflects the
finer grind the higher floor demands and affects every design equally,
as Fig.~3 of the paper also shows directly.

\begin{table}[H]
\centering
\caption{SEC of the six benchmark runs over the full post-warm-up
window, as in Table~I of the paper, and over the single-set-point
region after the 75~h floor step.}
\label{tab:sec75}
\begin{tabular}{llcc}
\toprule
Architecture & Objective & SEC, 50--218 h & SEC, 75--218 h \\
 & & (kWh/t) & (kWh/t) \\
\midrule
1-layer EMPC & energy & 19.58 & 19.92 \\
1-layer EMPC & profit & 19.55 & 19.93 \\
2-layer EMPC & energy & 19.64 & 19.98 \\
2-layer EMPC & profit & 19.31 & 19.60 \\
SSRTO+3PI    & energy & 19.65 & 19.94 \\
SSRTO+3PI    & profit & 19.58 & 19.90 \\
\bottomrule
\end{tabular}
\end{table}

The reviewer's broader point stands. If absolute economic percentages
were reported for an operating plant, the stability of $A$ over the
realised PSE distribution would have to be established with the
downstream circuits in the loop. That analysis is part of ongoing
plant-calibration work, and in follow-up work we plan to simplify the
protocol to a single set-point. Since this choice does not change our
main conclusions, we keep the reported results over the full region,
which keeps the tables and figures of the paper consistent.

\subsection*{Minor 1 -- Supporting values for the order-of-magnitude claim in III-B}

\begin{reviewquote}
Section III-B asserts that revenue exceeds energy cost by more than an
order of magnitude, but the supporting values ($A=71$~\$/t,
$B=0.06$~\$/kWh) appear only in Section V. A forward reference or
approximate magnitudes would help.
\end{reviewquote}

\textbf{Reply.} Agreed. The claim in Section III-B now carries a
forward reference, ``Revenue exceeds energy cost by more than an order
of magnitude (Section~V)'', pointing to Section~V where the values
$A=71$~\$/t and $B=0.06$~\$/kWh are given.

\subsection*{Minor 2 -- ``s.t.'' renders as $\sigma.\tau.$ in Eq.~(2)}

\begin{reviewquote}
In Eq.~(2), ``subject to'' appears to render as $\sigma.\tau.$ rather
than s.t., likely due to font substitution. Please verify.
\end{reviewquote}

\textbf{Reply.} This is also the typo reported by Reviewer 1, so both
reports are resolved by this fix.

\subsection*{Minor 3 -- Hyphenation breaks produce run-together words}

\begin{reviewquote}
Hyphenation breaks produce run-together words at several points (e.g.
``orehardness,'' ``singlestage,'' ``recedinghorizon'').
\end{reviewquote}

\textbf{Reply.} We could not find any of the run-together forms in
either the PDF or the Word file; all compound terms are correctly
hyphenated in both. The quoted forms are likely artifacts of joining
line-end hyphenated words while the paper was converted or annotated
during review.

\end{document}
